% from ega4-8-fr.tex
\oldpage[IV-3]{52}
par $\mathcal C'_0$-$\boldsymbol{gr}$ la catégorie des
$\mathcal C'_0$-groupes, il est vraisemblable que la catégorie des
$\mathcal C'$-groupes est équivalente à la sous-catégorie pleine de
$\mathbf{Pro}(\mathcal C'_0\text{-}\boldsymbol{gr})$ formée des «~groupes
proalgébriques essentiellement affines~», c'est-à-dire des systèmes
projectifs $\mathbf G=(G_\mu)_{\mu\in M}$, où les $G_\mu$ sont des groupes
algébriques sur $k$ et les morphismes de transition $G_\nu\to G_\mu$ sont
affines pour $\mu$ assez grand (ce que l'on peut exprimer aussi en disant que
$\mathbf G$ est extension d'un groupe algébrique par un pro-groupe algébrique
affine). La conjonction de ces deux conjectures équivaut d'ailleurs à la
suivante~: tout préschéma en groupes quasi-compacts sur $k$ est «~extension~»
d'un «~groupe algébrique~» (i.e. un préschéma en groupes de type fini sur $k$)
par un préschéma en groupes \emph{affine} sur $k$.

Les seuls pro-groupes algébriques rencontrés en pratique jusqu'à présent étant
en fait essentiellement affines, il y aura donc sans doute avantage à
substituer à l'étude des groupes pro-algébriques généraux (introduits et
étudiés par Serre [40]) celle des schémas en groupes quasi-compacts sur $k$,
dont la définition est conceptuellement plus simple.
\end{remark}

\subsection{Caractérisation d'un préschéma localement de présentation finie sur un autre, en termes du foncteur qu'il représente}
\label{subsection:IV.8.14-fr}

\begin{env}[8.14.1]
\label{IV.8.14.1-fr}
Étant donné un préschéma $S$, nous dirons encore, comme en (8.13.4), qu'un
système projectif filtrant $(X_\lambda,v_{\lambda\mu})$ de $S$-préschémas est
\emph{essentiellement affine} s'il existe $\alpha$ tel que
$v_{\alpha\lambda}$ soit un morphisme affine pour $\lambda\geqslant\alpha$.

L'énoncé suivant, qui sera surtout utile au chap.~V, précise (8.8.2, (i)) en
lui fournissant une réciproque~:
\end{env}

\begin{proposition}[8.14.2]
\label{IV.8.14.2-fr}
Soient $S$ un préschéma, $f:X\to S$ un morphisme. Pour tout $S$-préschéma
$T$, posons
\[
h_X(T)=\operatorname{Hom}_S(T,X)
\]
de sorte que $h_X$ est un foncteur contravariant de la catégorie
$(\mathbf{Sch})_{/S}$ des $S$-préschémas dans la catégorie
$\mathbf{Ens}$ des ensembles $(\mathbf 0_{\mathrm{III}},8.1.1)$, et $X$ un
objet représentant ce foncteur $(\mathbf 0_{\mathrm{III}},8.1.8)$. Les
conditions suivantes sont équivalentes~:
\begin{enumerate}[label=\emph{\alph*)}]
  \item $f$ est localement de présentation finie.
  \item Pour tout système projectif filtrant $(Z_\lambda)$ de
  $S$-préschémas, essentiellement affine (8.13.4) et formé de préschémas
  quasi-compacts et quasi-séparés, l'application canonique (8.13.1.1)
\begin{equation}
\label{IV.8.14.2.1-fr}
\varinjlim_\lambda h_X(Z_\lambda)
\longrightarrow h_X(\varprojlim_\lambda Z_\lambda) \tag{8.14.2.1}
\end{equation}
  est bijective.
  \item Pour tout système projectif filtrant $(Z_\lambda)$ de
  $S$-préschémas, tel que les $Z_\lambda$ soient des schémas affines,
  l'application (8.14.2.1) est bijective.
  \item[\rm c$'$)] Pour tout ouvert affine $U$ de $S$ et tout système
  projectif filtrant $(Z_\lambda)$ de $U$-préschémas, tel que les $Z_\lambda$
  soient des schémas affines, l'application (8.14.2.1) est bijective.
\end{enumerate}
\end{proposition}

Le fait que \emph{a)} implique \emph{b)} n'est autre que (8.13.1)~; il est
trivial que \emph{b)} implique \emph{c)} et que \emph{c)} implique
\emph{c$'$)}. Reste à voir que \emph{c$'$)} entraîne \emph{a)}, et comme la
propriété \emph{a)} est locale sur $S$, on peut se borner au cas où $S$ est
\emph{affine}.

